\documentclass[11pt]{article}
\usepackage[a4paper,margin=2.5cm]{geometry}
\usepackage[T1]{fontenc}
\usepackage{lmodern}
\usepackage{amsmath,amssymb,booktabs,longtable,array}
\usepackage{xcolor}
\usepackage{listings}
\usepackage[hidelinks]{hyperref}
\lstset{language=Matlab,basicstyle=\ttfamily\small,columns=fullflexible,keepspaces=true,frame=single,
  rulecolor=\color{black!30},breaklines=true,showstringspaces=false,keywordstyle=\color{black},
  commentstyle=\color{black!60}}
\newcommand{\prog}[1]{\texttt{#1}}
\newcommand{\tp}{^{\mathsf T}}
\newcommand{\He}{\mathrm{He}}
\newcommand{\K}{\mathcal{K}}
\newcommand{\pr}{\mathrm{P}}
\title{SaddlepointExpansions for MATLAB\\[4pt]
\large User manual}
\author{Simon A. Broda and Raymond Kan}
\date{October 2026}

\begin{document}
\sloppy
\maketitle

\tableofcontents

\section{Introduction}

SaddlepointExpansions for MATLAB computes saddlepoint expansions of any order for the
distribution of a random vector from its cumulant generating function alone. It implements the
expansions of Broda and Kan, ``Saddlepoint expansions of arbitrary order for multivariate
distribution functions'':
\begin{itemize}
\item the orthant probability $\pr(X\geq x)$ of a random vector $X$ of any dimension $d$, with lower
and mixed corners, and with the inclusion--exclusion step for saddlepoints that have negative
components;
\item the density of $X$ in any dimension;
\item the joint density and the joint distribution function of a vector of ratios
$R=(X_1/Y_1,\ldots,X_d/Y_d)$ with positive denominators, including ratios with a common
denominator;
\item in one dimension, the expansions of Daniels (density), Hauschildt and Robinson (distribution
function) and Rice (density of a ratio).
\end{itemize}
The approximation of order $M$ has relative error $O(n^{-M})$, where $n$ is the large parameter (a
sample size, a number of degrees of freedom).

The user writes the cumulant generating function $K$ as an ordinary MATLAB function. Its
derivatives of all orders are generated by Taylor-mode automatic differentiation (the class
\prog{tps}, truncated multivariate power series with overloaded arithmetic and elementary
functions), directly in the whitened coordinates in which the expansions are evaluated, so that no
derivative has to be written by hand and none is computed by finite differences. The expansions
themselves are computed by the MATLAB programs that accompany the paper (folder \prog{engine}),
which can also be called directly with closed-form derivatives.

The package is the MATLAB counterpart of the Julia package \prog{SaddlepointExpansions.jl} by
Simon Broda (\url{https://github.com/s-broda/SaddlepointExpansions.jl}). The two packages have
the same interface and agree to about $10^{-13}$ relative on all the test cases of
\S\,\ref{sec:valid}. The MATLAB package adds the expansions for several ratios
(\S\,\ref{sec:ratios}); the Julia package also runs in variable precision
(\S\,\ref{sec:precision}).

Theorem and corollary numbers refer to the version of the paper of October 2026: Theorem~1
(orthant probability), Theorem~2 (inclusion--exclusion), Theorem~3 (density), Theorem~4 and
Corollary~1 (densities of ratios), Corollaries~2 to~5 (Daniels, Hauschildt--Robinson, the
distribution function of a ratio, Rice).

\paragraph{Requirements.} MATLAB (the package was tested with R2025b) and the
Statistics and Machine Learning Toolbox, whose \prog{normcdf} and \prog{normpdf} are used by the
univariate engine. Nothing else is needed.

\paragraph{Installation.} Copy the folder anywhere and run \prog{spx\_setup}, which adds the folder
and its \prog{engine} subfolder to the path. \prog{help SaddlepointExpansions} lists the functions,
\prog{examples/demo\_spx.m} runs the examples of this manual in a few seconds and
\prog{tests/test\_spx.m} the tests of \S\,\ref{sec:valid} in about ten.

\section{A first example}\label{sec:first}

Let $X=(U_0+U_1,U_0+U_2)$ with $U_0,U_1,U_2$ independent unit exponential variables, the example
of Kolassa and Li (2010), and let $\bar X$ be the mean of $n=5$ independent copies of $X$. Then
\begin{lstlisting}
spx_setup
K = CGF(@(t) -log(1 - t(1)) - log(1 - t(2)) - log(1 - t(1) - t(2)), 2, ...
        'domain', @(t) t(1) < 1 && t(2) < 1 && t(1) + t(2) < 1);
Kn = sample_mean(K, 5);               % the cgf 5 K(t/5) of the mean
r = tailprob(Kn, [3 3], 'order', 4);  % P(mean1 >= 3, mean2 >= 3)
r.F'                                  % 0.02142343 0.02214365 0.02215413 0.02215193
r.tau                                 % the saddlepoint
s = density(Kn, [3 3], 'order', 4);   % the joint density of the mean at (3, 3)
\end{lstlisting}
The exact probability is $0.02215196$; the approximations of orders one to four have relative
errors $-3.3\times10^{-2}$, $-3.8\times10^{-4}$, $9.8\times10^{-5}$ and about $-1.4\times10^{-6}$. The
same cumulant generating function is in the library as
\prog{exponential\_sums\_cgf([1 0; 1 1; 0 1], 'n', 5)}.

\section{Conventions}\label{sec:conv}

\paragraph{The large parameter.} The large parameter $n$ is absorbed in the cumulant generating
function: the expansions are those of the vector whose cumulant generating function is supplied.
For the mean of $n$ copies one supplies $nK(t/n)$, which \prog{sample\_mean(K, n)} constructs; the
density returned is then that of the mean.

\paragraph{Orders.} Every expansion function takes \prog{'order', M} (default three) and returns the
approximations of orders $1,\ldots,M$ in a column vector. The approximation of order $M$ has
relative error $O(n^{-M})$. For a density it keeps the terms $m=0,\ldots,M-1$ of the series in
powers of $n^{-1}$, and for an orthant probability the terms $m=0,\ldots,2M-1$ of the series in
powers of $n^{-1/2}$. For densities, order one is Daniels's first-order approximation; for orthant
probabilities, the leading term $m=0$ is the approximation of Kolassa (2003).

\paragraph{Results.} The expansion functions return a structure with the approximations (\prog{F}
or \prog{f}), the saddlepoint (\prog{tau}, or \prog{shat} for ratios), a \prog{status} string
(\prog{'ok'} or a warning about the numerical evaluation) and \prog{info}, the diagnostics of the
underlying engine program. The diagnostics measure the numerical quality of the evaluation, such
as cancellation among the terms or the change with half the quadrature nodes; they do not measure
the asymptotic error of the expansion, which is what the order controls.

\paragraph{Vectors.} Points may be row or column vectors. The argument \prog{t} of a cumulant
generating function is always a $d\times1$ column vector.

\section{Cumulant generating functions}\label{sec:cgf}

\subsection{The class CGF}

\begin{lstlisting}
K = CGF(fun, d)
K = CGF(fun, d, 'domain', pred)
K = CGF(fun, d, 'Taylor', tfun)
\end{lstlisting}
\prog{fun} is a function of a $d\times1$ vector \prog{t} returning $K(t)$, also for $d=1$.
\prog{pred} is an optional predicate \prog{@(t) logical} describing the domain, and \prog{tfun} an
optional function that returns Taylor polynomials in closed form (\S\,\ref{sec:taylorfun}).

The function is evaluated on vectors of truncated power series to generate its derivatives, so it
must be written with operations that the class \prog{tps} supports:
\begin{itemize}
\item \prog{+}, \prog{-}, \prog{.*}, \prog{./}, \prog{.\^{}}, \prog{*} and \prog{/} by numbers,
products \prog{A*t} and \prog{t.'*A} with constant matrices, \prog{sum}, \prog{prod}, \prog{dot},
indexing \prog{t(j)}, \prog{t(1:k)}, concatenation, \prog{reshape}, \prog{diag};
\item \prog{exp}, \prog{expm1}, \prog{log}, \prog{log1p}, \prog{sqrt}, \prog{sin}, \prog{cos},
\prog{tan}, \prog{sinh}, \prog{cosh}, \prog{tanh}, \prog{atan}, \prog{abs}, and \prog{sinhc},
\prog{sincu} of the package;
\item comparisons (\prog{<}, \prog{<=}, \ldots), which compare the constant terms, so that code
with tests on the argument can be differentiated.
\end{itemize}
Functions of \prog{t} that call compiled numerical routines (\prog{det}, \prog{inv}, \prog{chol},
\prog{eig}, \prog{expm}, \ldots) cannot be differentiated this way; \S\,\ref{sec:own} shows how to
write determinants and inverses.

\paragraph{The domain.} Outside the domain, \prog{fun} may raise an error with identifier
\prog{'tps:domain'} (the logarithm and the square root of a nonpositive number do so
automatically) or return a non-finite value; both are treated as ``outside the domain'', and the
saddlepoint iteration steps back. When the formula is finite outside the true domain, as
$1/(1-t)$ is for $t>1$, the predicate \prog{'domain'} is needed. It is also cheaper than an error.

\subsection{The library}

\begin{longtable}{@{}>{\raggedright\arraybackslash}p{0.42\textwidth}>{\raggedright\arraybackslash}p{0.53\textwidth}@{}}
\toprule
constructor & distribution\\
\midrule
\endhead
\prog{normal\_cgf(mu, Sigma)} & $N(\mu,\Sigma)$ (the expansions are exact at order one)\\
\prog{gamma\_cgf(a, b)} & gamma with shape $a$ and scale $b$ (default one); \prog{gamma\_cgf(1, b)} is the exponential distribution with mean $b$\\
\prog{ncx2\_cgf(m, lam)} & noncentral chi-squared $\chi^2_m(\lambda)$\\
\prog{exponential\_sums\_cgf(A, 'n', n, 'shift', l)} & the mean of $n$ independent copies of $A\tp U+l$, $U$ a vector of independent unit exponential variables\\
\prog{quadratic\_forms\_cgf(Mats, 'Sigma', S, 'mu', mu, 'scale', c)} & the quadratic forms $z\tp A_iz$, $i=1,\ldots,d$, in $z\sim N(\mu,\Sigma)$, \prog{Mats = \{A\_1,...,A\_d\}}, summed over $c$ independent copies\\
\prog{wishart\_diag\_cgf(V, n)} & the diagonal of a Wishart matrix $W_d(n,V)$\\
\prog{common\_component\_cgf(m, lam)} & $Q_j=X_j+X_0$, $j=1,\ldots,d$, with independent $X_i\sim\chi^2_{m_i}(\lambda_i)$, \prog{m = [m\_1 ... m\_d m\_0]}\\
\prog{uniform\_cgf(a, b)} & uniform on $[a,b]$ (default $[-1,1]$)\\
\bottomrule
\end{longtable}

\noindent The cumulant generating function of quadratic forms is
\[
K(t)=-\frac12\log\det(I-2C)+\nu\tp C(I-2C)^{-1}\nu,\qquad C=\sum_{i=1}^{d}t_i\Sigma^{1/2}A_i\Sigma^{1/2},\qquad \nu=\Sigma^{-1/2}\mu,
\]
on the set where $I-2C$ is positive definite. Its Taylor polynomials are computed in closed form
by the trace recursion of the paper, which is much faster than automatic differentiation for
matrices of large order.

\subsection{Transformations}

\begin{longtable}{@{}>{\raggedright\arraybackslash}p{0.32\textwidth}>{\raggedright\arraybackslash}p{0.63\textwidth}@{}}
\toprule
function & cumulant generating function of\\
\midrule
\endhead
\prog{reflect(K, s)} & $(s_1X_1,\ldots,s_dX_d)$, $s_j=\pm1$; \prog{reflect(K)} is that of $-X$\\
\prog{marginal(K, idx)} & the components \prog{idx} of $X$\\
\prog{linear\_combination(K, B)} & $B\tp X$ for a $d\times k$ matrix $B$, that is $K(Bu)$\\
\prog{sample\_mean(K, n)} & the mean of $n$ independent copies, $nK(t/n)$\\
\prog{shift(K, c)} & $X+c$\\
\prog{K1 + K2} & the sum of two independent vectors of the same dimension\\
\bottomrule
\end{longtable}

\noindent The transformations keep closed-form Taylor polynomials: a linear combination of
quadratic forms, for example, is still differentiated by the trace recursion.

\subsection{Writing your own}\label{sec:own}

Removable singularities are the one trap: $\sinh(t)/t$ is $0/0$ at $t=0$, where the saddlepoint
iteration starts. Use \prog{sinhc} and \prog{sincu}, which switch to the power series there; the
library defines the uniform distribution as \prog{log(sinhc(t))}. Two patterns:
\begin{lstlisting}
% independent parts, one per component, plus a common part
K = CGF(@(t) -sum(m(:)/2.*log(1 - 2*t)) - m0/2*log(1 - 2*sum(t)), d, ...
        'domain', @(t) all(t < 1/2) && sum(t) < 1/2);

% a log-determinant: a Cholesky factorization written in plain arithmetic
function v = logdet_cgf(t, B1, B2)
n = size(B1, 1); Q = eye(n) - 2*(t(1)*B1 + t(2)*B2);
L = Q(1, 1)*zeros(n);                 % a series array of the right class
for j = 1:n
   s = Q(j, j); for k = 1:j-1, s = s - L(j, k)^2; end
   L(j, j) = sqrt(s);
   for i = j+1:n
      w = Q(i, j); for k = 1:j-1, w = w - L(i, k)*L(j, k); end
      L(i, j) = w/L(j, j);
   end
end
v = -sum(log(diag(L)));
end
\end{lstlisting}
The test file \prog{tests/test\_spx.m} differentiates the complete cumulant generating function of a
pair of quadratic forms written this way and checks it against the trace recursion. For quadratic
forms in normal variables, \prog{quadratic\_forms\_cgf} is faster.

\subsection{Closed-form Taylor polynomials}\label{sec:taylorfun}

The option \prog{'Taylor', tfun} supplies a function \prog{tfun(t, N, M)} that returns the Taylor
polynomial of degree $N$ in $v$ of $K(t+Mv)$ as a scalar \prog{tps} in \prog{size(M, 2)} variables
(\prog{M = []} stands for the identity). It is then used instead of automatic differentiation;
\prog{fun} is still used for the values of $K$. The polynomial is built from its coefficients with
\prog{tps(C, [1 1], k, N)}, where \prog{C} holds the coefficients of the monomials in the order of
\prog{tps.tables(k, N).mi} (\S\,\ref{sec:tps}). \prog{quadratic\_forms\_cgf} is an example.

\section{Orthant probabilities}\label{sec:tail}

\begin{lstlisting}
r = tailprob(K, x, 'order', M)                  % P(X >= x)
r = tailprob(K, x, 'order', M, 'lower', true)   % P(X <= x)
r = tailprob(K, x, 'order', M, 'signs', s)      % e.g. s = [1 -1]: P(X1 >= x1, X2 <= x2)
\end{lstlisting}
\prog{r.F} holds the approximations of orders $1,\ldots,M$ and \prog{r.tau} the saddlepoint,
$K'(\hat\tau)=x$. Lower and mixed corners are computed as upper orthant probabilities of the
reflected vector.

\paragraph{Saddlepoints with negative components.} The expansion of Theorem~1 applies to
saddlepoints in the closed positive orthant $\hat\tau\geq0$. When some components are negative,
\prog{tailprob} writes the probability by the inclusion--exclusion identity of Theorem~2 as a
signed sum of orthant probabilities of marginals with some components negated, each with a
saddlepoint in its closed positive orthant, recursively, and expands every term
(\S\,\ref{sec:background}). This keeps the relative error $O(n^{-M})$; a direct evaluation of the
formula at such a saddlepoint would have the wrong exponential order. The saddlepoint reported is
that of the original problem.

\paragraph{Dimensions.} In one dimension the expansion is that of Hauschildt and Robinson
(Corollary~3), computed directly for the upper tail. In two dimensions the integrals of the
expansion are computed by Gauss--Legendre quadrature (\prog{spa\_orthant2}, \prog{spa\_bi}), in
three by Gauss quadrature on a conditioned grid (\prog{spa\_tri}), and in four or more by randomized
quasi-Monte Carlo after Genz's sequential conditioning (\prog{spa\_mv}); \prog{r.info.err} is then an
error estimate (in four or more dimensions a standard error). A trivariate probability takes about
a second at order three and tens of seconds at order four or five, mostly in the inclusion--exclusion
terms.

\paragraph{Options.} Name-value pairs other than \prog{'order'}, \prog{'lower'} and \prog{'signs'}
are passed to the engine: \prog{'nG'} (Gauss--Legendre nodes per coordinate), \prog{'qtol'}
(quadrature tolerance), \prog{'nGmax'}, \prog{'method'} (\prog{'quad'} or \prog{'qmc'}),
\prog{'nqmc'}, \prog{'nshift'}, \prog{'seed'} and, in two dimensions, \prog{'logscale', c}, which
returns the probability scaled by $e^{-c}$ for probabilities below the range of double precision.
See the help of \prog{spa\_bi}, \prog{spa\_tri} and \prog{spa\_mv}.

\section{Densities}\label{sec:density}

\begin{lstlisting}
r = density(K, x, 'order', M)
\end{lstlisting}
\prog{r.f} holds the approximations of orders $1,\ldots,M$ of the density of $X$ at $x$
(Theorem~3; Daniels's expansion, Corollary~2, in one dimension). The expansion is a finite sum of
Hermite polynomials at zero and needs no integration in any dimension. In three or more dimensions
the terms of high orders can cancel by many orders of magnitude; \prog{r.info.kappa\_pdf} reports
the cancellation (\S\,\ref{sec:precision}).

\section{Ratios}\label{sec:ratios}

\subsection{Usage}

Let $R=(X_1/Y_1,\ldots,X_d/Y_d)$ with $Y_j>0$. The cumulant generating function \prog{K} is that of
the vector $(X_1,\ldots,X_d,Y_1,\ldots,Y_d)$, of dimension $2d$, numerators first. For ratios with a
common denominator $Y$ it may instead be that of $(X_1,\ldots,X_d,Y)$, of dimension $d+1$. The
number of ratios is the length of the point \prog{r0}.
\begin{lstlisting}
r = ratio_density(K, r0, 'order', M)                   % density of R at r0
r = ratio_tailprob(K, r0, 'order', M)                  % P(R >= r0)
r = ratio_tailprob(K, r0, 'order', M, 'lower', true)   % P(R <= r0)
r = ratio_tailprob(K, r0, 'order', M, 'signs', s)      % a mixed corner
\end{lstlisting}
\prog{r.f} or \prog{r.F} holds the approximations of orders $1,\ldots,M$, and \prog{r.shat} the
saddlepoint of $K_W$ (\S\,\ref{sec:bgratio}). For one ratio the density is Rice's expansion
(Corollary~5), and for several the expansion of Theorem~4 and Corollary~1, whose term of order one
is the approximation of Butler and Paolella (1998); the distribution function is the orthant
probability of \S\,\ref{sec:tail} for $W=X-r_0\circ Y$ (Corollary~4 for one ratio), with all its
options.

\subsection{Examples}

\paragraph{Sample autocorrelations.} With $z\sim N(0,I_N)$, $M=I_N-1_N1_N\tp/N$ and
$(A_j)_{ts}=1/2$ for $|t-s|=j$, the sample autocorrelations at lags one and two are
$R_j=z\tp MA_jMz/z\tp Mz$, ratios with the common denominator $z\tp Mz$. For $N=50$ at
$r_0=(-1,-1)/\surd50$:
\begin{lstlisting}
N = 50; M = eye(N) - ones(N)/N;
A = @(j) (diag(ones(N-j, 1), j) + diag(ones(N-j, 1), -j))/2;
K = quadratic_forms_cgf({M*A(1)*M, M*A(2)*M, M});      % dimension d + 1 = 3
r0 = [-1 -1]/sqrt(N);
f = ratio_density(K, r0, 'order', 5).f
F = ratio_tailprob(K, r0, 'order', 5, 'lower', true).F
\end{lstlisting}
The relative errors of orders one to five, against the exact values of the paper, are
\begin{center}
\begin{tabular}{lrrrrr}
\toprule
order & 1 & 2 & 3 & 4 & 5\\
\midrule
density & $-1.4\times10^{-3}$ & $7.3\times10^{-5}$ & $-1.0\times10^{-5}$ & $2.2\times10^{-6}$ & $-6.4\times10^{-7}$\\
distribution function & $-3.7\times10^{-2}$ & $1.5\times10^{-4}$ & $1.4\times10^{-5}$ & $5.8\times10^{-7}$ & $-1.6\times10^{-7}$\\
\bottomrule
\end{tabular}
\end{center}
These are the values of the corresponding row of the table of sample autocorrelations in the paper,
which were computed by a dedicated program.

\paragraph{Gamma ratios with a common denominator.} If $X_1,\ldots,X_d,Y$ are independent with
$X_j\sim\mathrm{Gamma}(a_j)$ and $Y\sim\mathrm{Gamma}(b)$, the ratios $R_j=X_j/Y$ have the inverted
Dirichlet density
\[
f_R(r)=\frac{\Gamma(a_1+\cdots+a_d+b)}{\Gamma(b)\prod_{j=1}^{d}\Gamma(a_j)}\,
\frac{\prod_{j=1}^{d}r_j^{a_j-1}}{(1+r_1+\cdots+r_d)^{a_1+\cdots+a_d+b}}.
\]
\begin{lstlisting}
a = [3 4 5]; b = 6;
K = CGF(@(t) -sum(a(:).*log(1 - t(1:3))) - b*log(1 - t(4)), 4, 'domain', @(t) all(t < 1));
f = ratio_density(K, [0.6 0.5 0.9], 'order', 5).f
\end{lstlisting}
The relative errors of orders one to five are $7.7\times10^{-2}$, $-3.1\times10^{-3}$,
$-1.0\times10^{-4}$, $1.7\times10^{-5}$ and $3.6\times10^{-6}$. Multiplying the shapes by $n$ shows
the rate $O(n^{-M})$; for $d=2$ at $r_0=(0.6,0.5)$:
\begin{center}
\begin{tabular}{rrrrr}
\toprule
$n$ & order 1 & order 2 & order 3 & order 4\\
\midrule
1 & $5.8\times10^{-2}$ & $-1.8\times10^{-3}$ & $-1.2\times10^{-4}$ & $1.3\times10^{-5}$\\
2 & $2.8\times10^{-2}$ & $-4.2\times10^{-4}$ & $-1.6\times10^{-5}$ & $6.6\times10^{-7}$\\
4 & $1.4\times10^{-2}$ & $-1.0\times10^{-4}$ & $-2.0\times10^{-6}$ & $3.7\times10^{-8}$\\
8 & $7.0\times10^{-3}$ & $-2.5\times10^{-5}$ & $-2.5\times10^{-7}$ & $2.2\times10^{-9}$\\
\bottomrule
\end{tabular}
\end{center}
Each doubling of $n$ divides the error of order $M$ by about $2^M$.

\subsection{Cost}

The density of several ratios needs one Taylor polynomial of $\K(s,-r\circ s+\tau)$ in the $2d$
variables $(u,\tau)$ of \S\,\ref{sec:bgratio}, of degree $\max(2M,2M-2+d)$. For two or three
ratios every order up to five takes one or two seconds, also for quadratic forms in $N=50$ variables;
for four ratios order four takes several seconds. The distribution function costs what an orthant
probability of dimension $d$ costs (\S\,\ref{sec:tail}).

\section{Derivatives, saddlepoints and the Taylor class}\label{sec:tps}

\begin{lstlisting}
D = derivatives(K, t, N)       % all derivatives of orders <= N at t
[tau, info] = saddlepoint(K, x)
p = taylor(K, t, N, M)         % Taylor polynomial of K(t + M v) in v, degree N
\end{lstlisting}
\prog{D.nu} holds the multi-indices $\nu$, one per row in graded order, and \prog{D.value} the
derivatives $K^{(\nu)}(t)$, so that \prog{D.value(1)} is $K(t)$. The saddlepoint is computed by
Newton's method with backtracking on the convex function $K(\tau)-\tau\tp x$ (\prog{spa\_saddle}).

The class \prog{tps} represents arrays of polynomials in $k$ variables truncated at total degree
$N$:
\begin{lstlisting}
v = tps.vars(k, N)       % k x 1 array of the variables
p = exp(v(1)) * log(1 + v(2))
coef(p, [1 2])           % nu! times the coefficient of v^nu, the derivative at 0
double(p)                % the constant term
T = tps.tables(k, N);    % T.mi: the multi-indices in graded order
\end{lstlisting}
A function $f$ of a series $p=p_0+q$, $q$ without constant term, is evaluated by composition,
$f(p)=\sum_{j=0}^{N}f^{(j)}(p_0)q^j/j!$; where $f$ is not analytic at $p_0$ an error with identifier
\prog{'tps:domain'} is raised. The coefficients are stored in \prog{p.C}, one column per element of
the array.

\section{The expansions}\label{sec:background}

This section states what the package computes. The paper gives the derivations, the conditions and
the error bounds; here $n$ is absorbed in $K$.

\subsection{Whitened coefficients}

Let $\hat\tau$ be the saddlepoint, $K'(\hat\tau)=x$, and $K''(\hat\tau)=LL\tp$ the Cholesky
factorization. All terms are computed for the whitened vector $Y=L^{-1}X$, whose cumulant generating
function $K_Y(u)=K(L^{-\mathsf T}u)$ has its saddlepoint at $\hat u=L\tp\hat\tau$ and Hessian $I$
there. Let $\theta_\nu=K_Y^{(\nu)}(\hat u)/\nu!$, $|\nu|\geq3$, be its Taylor coefficients and define
$c_{r,\nu}$ by
\[
c_{1,\nu}=\theta_\nu,\qquad
c_{r,\nu}=\frac1r\sum_{k=3}^{|\nu|-3(r-1)}\ \sum_{|\lambda|=k,\ \lambda\leq\nu}\theta_\lambda\,c_{r-1,\nu-\lambda},\qquad r\geq2,
\]
so that $c_{r,\nu}$ is the coefficient of $u^\nu$ in $\Theta(u)^r/r!$, where
\[
\Theta(u)=\sum_{k=3}^{\infty}\ \sum_{|\nu|=k}\theta_\nu u^\nu.
\]
The multivariate Hermite polynomials at zero factor into univariate ones,
\[
\He_\nu(0)=\prod_{j=1}^{d}\He_{\nu_j}(0),\qquad \He_{2i}(0)=(-1)^i\prod_{j=1}^{i}(2j-1),\qquad \He_{2i+1}(0)=0.
\]
The package obtains the $\theta_\nu$ by differentiating the composed function $K(\hat\tau+L^{-\mathsf T}u)$
directly. In the original coordinates the terms cancel much more when the components of $X$ are
strongly correlated under the tilted distribution, and the substitution of $L^{-\mathsf T}$ into the
Taylor polynomial of $K$ can lose many digits at high degree.

\subsection{Densities}

\[
f(x)\approx\frac{\exp\{K(\hat\tau)-\hat\tau\tp x\}}{(2\pi)^{d/2}|K''(\hat\tau)|^{1/2}}\sum_{m=0}^{M-1}a_m,\qquad
a_0=1,\qquad a_m=\sum_{r=1}^{2m}\ \sum_{|\nu|=2m+2r}c_{r,\nu}\He_\nu(0).
\]
The coefficient $a_m$ involves the $\theta_\nu$ with $|\nu|\leq2m+2$.

\subsection{Orthant probabilities}

For $\hat\tau\geq0$,
\[
\pr(X\geq x)\approx\exp\{K(\hat\tau)-\hat\tau\tp x+|\hat u|^2/2\}\sum_{m=0}^{2M-1}T_m,\qquad
T_0=\tilde I_0,\qquad T_m=\sum_{r=1}^{m}\ \sum_{|\nu|=m+2r}c_{r,\nu}\tilde I_\nu,
\]
with
\[
\tilde I_\nu=\int_{\{v\in\mathbb{R}^d:\ Lv\geq0\}}\ \prod_{j=1}^{d}\He_{\nu_j}(v_j)\,\phi_d(v+\hat u)\,dv,
\]
$\phi_d$ the standard normal density in $d$ dimensions. The term $T_m$ is of order $n^{-m/2}$.
When $\hat\tau$ has negative components, let $\mathcal N=\{j:\hat\tau_j<0\}$. Then
\[
\pr(X\geq x)=\sum_{T\subseteq\mathcal N}(-1)^{|T|}\,\pr\bigl(\tilde X_{\mathcal N^c\cup T}\geq\tilde x_{\mathcal N^c\cup T}\bigr),
\]
where $\tilde X$ and $\tilde x$ are $X$ and $x$ with the components in $T$ negated and the subscript
selects components; the term with $\mathcal N^c\cup T$ empty is one. Each term has a saddlepoint in
its closed positive orthant, or is decomposed again.

\subsection{Ratios}\label{sec:bgratio}

Let $\K(s,t)$, $s,t\in\mathbb{R}^d$, be the joint cumulant generating function of $(X,Y)$,
$W_j=X_j-r_jY_j$, and $r\circ s=(r_1s_1,\ldots,r_ds_d)$. The cumulant generating function of $W$ is
$K_W(s)=\K(s,-r\circ s)$, and $\pr(R\leq r)=\pr(W\leq0)$ is the orthant probability of $-W$ at
zero. For the density, let
\[
\K_B=\frac{\partial^{|B|}\K}{\prod_{j\in B}\partial t_j},\qquad B\subseteq\{1,\ldots,d\},
\]
and let $\mathcal P_{d,j}$ be the set of partitions of $\{1,\ldots,d\}$ into $j$ blocks and
\[
g_k(s)=\sum_{\pi\in\mathcal P_{d,d-k}}\ \prod_{B\in\pi}\K_B(s,-r\circ s),\qquad k=0,\ldots,d-1.
\]
With $\hat s$ the saddlepoint of $K_W$,
$K_W'(\hat s)=0$, and $K_W''(\hat s)=LL\tp$,
\[
f_R(r)\approx\frac{\exp\{K_W(\hat s)\}}{(2\pi)^{d/2}|K_W''(\hat s)|^{1/2}}\sum_{m=0}^{M-1}\tilde a_m,\qquad
\tilde a_m=\sum_{k=0}^{\min(m,d-1)}b_{m-k}(g_k),
\]
\[
b_m(h)=\sum_{r=0}^{2m}\ \sum_{j=3r}^{2m+2r}\ \sum_{|\nu|=j}\ \sum_{|\mu|=2m+2r-j}c_{r,\nu}\,\gamma_\mu\,\He_{\nu+\mu}(0),\qquad c_{0,0}=1,
\]
where $\theta_\nu$, and hence $c_{r,\nu}$, are the whitened coefficients of $K_W$ at $\hat s$ and
$\gamma_\mu$ the Taylor coefficients of $h(\hat s+L^{-\mathsf T}u)$ in $u$. The leading coefficient
is $\tilde a_0=g_0(\hat s)=\prod_{j=1}^{d}\K_{\{j\}}(\hat s,-r\circ\hat s)$. All the coefficients come
from one Taylor polynomial of
\[
\K\bigl(\hat s+L^{-\mathsf T}u,\ -r\circ(\hat s+L^{-\mathsf T}u)+\tau\bigr)
\]
in $(u,\tau)$: the coefficient of $\prod_{j\in B}\tau_j$ is the Taylor polynomial of $\K_B$ in $u$.
The sums $b_m$ are evaluated as $\He$-weighted sums of the coefficients of the products
$\Theta(u)^r\,h(u)/r!$.

\section{Accuracy and precision}\label{sec:precision}

The computations are in double precision. Rounding errors are about $10^{-14}$ relative in the
tests, far below the asymptotic errors of the expansions at the orders used in practice ($10^{-3}$ to
$10^{-9}$), so double precision suffices for nearly all purposes. It matters in three situations:
\begin{itemize}
\item \emph{Cancellation in three or more dimensions.} At high orders the terms can cancel by many
orders of magnitude. The engine reports this (\prog{kappa} in \prog{info}) and, when the
cancellation exceeds $10^6$, tries to repeat the computation in variable precision. The Taylor class
works in double precision only, so in this package that refinement is not possible: the
double-precision values are kept and \prog{status} reports \prog{'refinement failed'}. The Julia
package repeats the computation in \prog{BigFloat}.
\item \emph{Probabilities below about $10^{-308}$.} In two dimensions use \prog{'logscale'}
(\S\,\ref{sec:tail}).
\item \emph{Tables of very small relative errors}, at large $n$ and high order, need exact values and
expansions beyond 15 digits. The engine programs accept derivatives in variable precision
(\prog{vpa}, Symbolic Math Toolbox) when they are supplied in closed form (\S\,\ref{sec:engine}).
\end{itemize}

\section{The engine}\label{sec:engine}

The folder \prog{engine} contains the programs of the paper. They take the derivatives of $K$ from a
function \prog{Kfun(nu, t)}, the derivative of order \prog{nu} ($1\times d$) at \prog{t}
($1\times d$), \prog{NaN} outside the domain, and can be used directly when closed-form derivatives
are available:
\begin{center}
\begin{tabular}{@{}ll@{}}
\toprule
\prog{spa\_saddle} & saddlepoint by Newton's method with backtracking\\
\prog{spa\_uni} & Daniels and Hauschildt--Robinson expansions from $K^{(k)}(\hat\tau)$\\
\prog{spa\_ratio\_pdf} & Rice's expansion of the density of a ratio\\
\prog{spa\_bi}, \prog{spa\_orthant2} & bivariate density and orthant probability\\
\prog{spa\_tri}, \prog{spa\_mv}, \prog{spa\_orthant} & three and more dimensions\\
\prog{qf\_Kmulti} & derivatives of the cgf of quadratic forms\\
\prog{ncx2\_spa}, \prog{ncf\_spa}, \prog{dncf\_spa} & noncentral $\chi^2$, $F$ and doubly noncentral $F$\\
\prog{spa\_uni\_hp}, \prog{spa\_bi\_hp} & variable precision\\
\bottomrule
\end{tabular}
\end{center}
The class \prog{CGF} builds \prog{Kfun} and the whitened derivatives (\prog{opts.Kwhite}) from the
Taylor polynomials. See the help of each program.

\section{Validation}\label{sec:valid}

\prog{tests/test\_spx.m} runs the following checks and returns \prog{true} if all pass
(\prog{test\_spx(true)} prints every result):
\begin{itemize}
\item derivatives of the Taylor class against closed forms and finite differences, and the
derivatives of a pair of quadratic forms by automatic differentiation through a generic Cholesky
factorization against the trace recursion;
\item 103 values of the Julia package (\prog{tests/julia\_reference.csv}, written by
\prog{tests/julia\_reference.jl}): orthant probabilities, lower and mixed corners, densities and
single ratios for sums of exponentials in two and three dimensions, gamma, noncentral chi-squared,
uniform, common-component, quadratic-form, Wishart and normal distributions. All agree to about
$10^{-13}$ relative;
\item the expansion for several ratios at $d=1$ against Rice's expansion ($10^{-14}$), and for two and
three ratios against the inverted Dirichlet density and distribution function and the product of
two independent beta-prime ratios, including the rate $O(n^{-4})$ of the density of order four.
\end{itemize}
In addition, for all 24 points of the table of sample autocorrelations of the paper, the density and
distribution function of two ratios computed by the package agree with those of the dedicated
program of the paper (\prog{acf\_spa}) to $1.5\times10^{-12}$ at every order from one to five.

\section{Function reference}

\begin{longtable}{@{}>{\raggedright\arraybackslash}p{0.42\textwidth}>{\raggedright\arraybackslash}p{0.53\textwidth}@{}}
\toprule
call & result\\
\midrule
\endhead
\prog{K = CGF(fun, d, 'domain', pred, 'Taylor', tfun)} & a cumulant generating function\\
\prog{r = tailprob(K, x, 'order', M, 'lower', tf, 'signs', s, ...)} & orthant probabilities; \prog{r.F}, \prog{r.tau}, \prog{r.status}, \prog{r.info}\\
\prog{r = density(K, x, 'order', M, ...)} & densities; \prog{r.f}, \prog{r.tau}, \prog{r.status}, \prog{r.info}\\
\prog{r = ratio\_density(K, r0, 'order', M)} & densities of $d=$ \prog{numel(r0)} ratios; \prog{r.f}, \prog{r.shat}, \prog{r.info} (for $d\geq2$: the terms \prog{a}, the log prefactor \prog{logC})\\
\prog{r = ratio\_tailprob(K, r0, 'order', M, 'lower', tf, 'signs', s, ...)} & distribution functions of ratios; \prog{r.F}, \prog{r.shat}, \prog{r.status}, \prog{r.info}\\
\prog{[tau, info] = saddlepoint(K, x)} & the saddlepoint\\
\prog{D = derivatives(K, t, N)} & \prog{D.nu}, \prog{D.value}\\
\prog{p = taylor(K, t, N, M)} & Taylor polynomial of $K(t+Mv)$\\
\prog{v = eval(K, t)} & $K(t)$\\
\prog{reflect}, \prog{marginal}, \prog{linear\_combination}, \prog{sample\_mean}, \prog{shift}, \prog{+} & transformations (\S\,\ref{sec:cgf})\\
\prog{tps.vars(k, N)}, \prog{tps.const(c, k, N)}, \prog{coef(p, nu)}, \prog{double(p)}, \prog{tps.tables(k, N)} & the Taylor class\\
\prog{sinhc(x)}, \prog{sincu(x)} & $\sinh(x)/x$ and $\sin(x)/x$\\
\prog{spx\_setup} & adds the package to the path\\
\bottomrule
\end{longtable}

\section*{References}
\begin{description}\setlength{\itemsep}{2pt}
\item Broda, S. A. and Kan, R. Saddlepoint expansions of arbitrary order for multivariate
distribution functions. Working paper, 2026.
\item Butler, R. W. (2007). \emph{Saddlepoint Approximations with Applications}. Cambridge
University Press.
\item Butler, R. W. and Paolella, M. S. (1998). Approximate distributions for the various serial
correlograms. \emph{Bernoulli} 4, 497--518.
\item Daniels, H. E. (1954). Saddlepoint approximations in statistics. \emph{Ann. Math. Statist.} 25,
631--650.
\item Hauschildt, R. (1969). Saddlepoint approximations to distribution functions. Master's thesis,
University of British Columbia, Vancouver.
\item Kolassa, J. E. (2003). Multivariate saddlepoint tail probability approximations. \emph{Ann.
Statist.} 31, 274--286.
\item Kolassa, J. E. and Li, J. (2010). Multivariate saddlepoint approximations in tail probability and
conditional inference. \emph{Bernoulli} 16, 1191--1207.
\end{description}

\end{document}
